You can view my CV in \href{cv.md}{PDF} file or see my
\href{http://www.linkedin.com/in/ching-chi-chou-31977b185}{LinkedIn}.

\subsection{\#\#\# Contact}\label{contact}

Email: ching-chi.chou@epfl.ch

Phone: +41 76 221 89 17

\subsection{\#\#\# Education}\label{education}

\textbf{EPFL (Swiss Federal Institute of Technology, Lausanne)}\\
M.S. in Digital Humanities (09.2024 - Present)

\textbf{National Taiwan University}\\
B.S. in Geography - Geoinformation Track (09.2019 - 06.2023)

\begin{itemize}
\item
  Specialization Program in Geographic Data Science
\item
  Course: Machine Learning, Programming Design, Operating Systems,
  Spatial Database Systems, Web Programming, Big Data Analytics, Network
  Data Analysis and Models, UAS Mapping and Image Interpretation,
  Geospatial Simulation, Remote Sensing, Spatiotemporal Data
  Visualization, Cartography and GIS
\end{itemize}

\subsection{\#\#\# Experience}\label{experience}

\textbf{Center for GIS, RCHSS, Academia Sinica}\\
Research Intern (07.2022 - 06.2023)

\begin{itemize}
\item
  Utilized spatiotemporal statistics and telecom mobility metrics,
  including centrality, structural equivalence, and footprint
  similarity, to assess policy intervention effectiveness and predict
  hotspots during the COVID-19 outbreak in 2021.
\item
  Simulated spatiotemporal epidemic progression by the SEPIAQR
  infectious disease model and mobility data, with proposed resource
  allocation scenarios achieving 20\% fewer confirmed cases with
  optimized testing resource usage.
\end{itemize}

\textbf{H2U Corporation}\\
Digital Product Manager (09.2022 - . 12.2023)

\begin{itemize}
\item
  Transformed digital products into \textbf{subscription-based} services
  by launching a gamified hiking route-tracking service with 100K+
  monthly active users, combining offline maps with local cultural and
  ecological features.
\item
  Collaborated with the Ministry of Digital Affairs to promote a public
  sports data platform, and won over NT\$500K in a sports data
  hackathon.
\end{itemize}

Data Analyst Intern (07.2022 - 09.2022)

\begin{itemize}
\item
  Developed a spatiotemporal data analysis tool, \textbf{Hiking Tracks
  Explorer}, featuring a hiking heatmap and trail popularity prediction
  using user GPS data and web-crawled GPS data.
\item
  Built online dashboards using Google Analytics and Looker Studio,
  providing insights to optimize advertising strategy.
\end{itemize}

\subsection{\#\#\# Publications}\label{publications}

\textbf{Peer-reviewed Publications}

{[}1{]} Chan, T. C., \textbf{Chou, C. C.}, Chu, Y. C., Tang, J. H.,
Chen, L. C., Lin, H. H., \ldots{} \& Chen, R. C. (2022).
\emph{Effectiveness of controlling COVID-19 epidemic by implementing
soft lockdown policy and extensive community screening in Taiwan.}
Scientific Reports, 12(1), 1-11.
\href{https://www.nature.com/articles/s41598-022-16011-x}{Link}

{[}2{]} \textbf{Chou, C. C.}, Lin, H. H., Chen, K. J., Chen, R. C.,
Chan, T. C. (2022). \emph{Comparing performance between static and
dynamic populations applied to COVID-19 hotspot prediction.} Taiwan
Journal Public Health, 41(6), 611-626.
\href{http://doi.org/10.6288/TJPH.202212_41(6).111074}{Link}

{[}3{]} Shih, T. L., \textbf{Chou, C. C.}, Chen, Y. C. (2018).
\emph{Living or Leaving? The Alteration and Perception of Place Names of
Kavalan in Yilan, Taiwan.} Bulletin of The Geographical Society of
China, 61, 51-66. \href{https://bit.ly/3JTOtxO}{Link}

\textbf{Conference Publication}

{[}4{]} \textbf{Chou, C. C.}, Lin, H. H., Chen, K. J., Chen, R. C.,
Chan, T. C. (2022). \emph{Predicting the hotspots of COVID-19 outbreak
by Telecom data - Take Taipei City and New Taipei City as an example.}
Taiwan Geographic Information Conference 2022, 89-96.
\href{https://drive.google.com/file/d/1KBwXkAsi8dmpPls-F3Lp8rqz5oc2JKWW/view}{Link}
\emph{(Best Paper Award)}

\textbf{Non-peer-reviewed Report}

{[}5{]} \textbf{Chou, C. C.} (2023). \emph{Exploring the spatial
allocation of testing stations using supply-demand relations: A Case
Study of New Taipei City.} College Student Research Grant, National
Science and Technology Council.

\subsection{\#\#\# Awards \& Honors}\label{awards-honors}

\textbf{College Student Research Grant, National Science and Technology
Council} (07.2022 - 02.2023)

\begin{itemize}
\item
  Project Title: \emph{Exploring the spatial allocation of testing
  stations using supply-demand relations: A Case Study of New Taipei
  City}
\item
  Total amount of grant: NT\$48,000
\end{itemize}

\textbf{Honorable Mention, Move to Earn Sport Data Hackathon} (02.2023)

\begin{itemize}
\item
  Represented H2U Corporation, awarded by Ministry of Digital Affairs
\item
  Prevailed over 40+ companies to win NT\$500,000
\end{itemize}

\textbf{Best Paper Award, Taiwan Geographic Information Conference}
(07.2022)

\textbf{Selected Award, Keelung City Hackathon for Geo-Science Strategy}
(01.2022)

\textbf{8th in 3000m Steeplechase, National Intercollegiate Athletic
Games} (11.2020)

\textbf{Silver Medal, National Geography Olympiad} 10.2018

\subsection{\#\#\# Extracurricular
Activities}\label{extracurricular-activities}

\textbf{International Companions for Learning} (03.2023 - 06.2023)

\begin{itemize}
\tightlist
\item
  Cooperated with an Indonesian companion to help elementary students in
  remote areas.
\end{itemize}

\textbf{Coach in EMBA Running Club, National Chengchi University}
(Sep.~2022 - Dec.~2023 )

\begin{itemize}
\tightlist
\item
  Arranged customized workout schedules for 40+ club members and won
  multiple road running events in Taiwan.
\end{itemize}

\textbf{Compulsory Military Service, Taiwan Army} (Dec.~2023 -
Apr.~2024)

\subsection{\#\#\# Additional Skills}\label{additional-skills}

\begin{itemize}
\item
  \textbf{Data Science:} Python/Pytorch, R, SQL
\item
  \textbf{Software Development:} C++, React.js, Node.js
\item
  \textbf{GIS-related:} QGIS, ArcGIS, NetLogo, Stella Architect,
  Metashape,Mapbox-gl.js, Leaflet.js, Google Earth Engine
\item
  \textbf{Tools:} Git, LaTex, Tableau, Power BI, Looker Studio, Google
  Analytics, Figma
\end{itemize}
